\documentclass[10pt,a4paper]{amsart}
\usepackage[margin=19mm]{geometry}
\usepackage{amsmath,amssymb,mathtools}
\usepackage[scr=rsfs,cal=euler]{mathalfa}
\usepackage{xfrac}
\usepackage[unicode,hidelinks,bookmarks=false]{hyperref}
\newcommand{\ie}{i.e.\ }
\newcommand{\cf}{cf.\ }
\newcommand{\resp}{respectively\ }
\newcommand{\Subsection}{\subsection}
\providecommand{\nobreakdash}{\nobreak\hbox{-}}
\setlength{\emergencystretch}{2em}
\multlinegap0pt
\usepackage{enumerate}
\usepackage{amssymb}
\usepackage{bm}
\usepackage{cancel}
\usepackage{tikz}
\newtheorem{thm}{Theorem}
\newtheorem{lem}{Lemma}
\title{Existence and nonexistence results for a nonlocal isoperimetric problem on $\mathbb{H}^n$}
\author{Haizhong Li, Bo Yang}
\date{Source: arXiv:2512.12621v2 (CC BY 4.0)}
\begin{document}
\maketitle

\noindent\emph{Source and licence.} Existence and nonexistence results for a nonlocal isoperimetric problem on $\mathbb{H}^n$. Version arXiv:2512.12621v2 (CC BY 4.0), distributed by arXiv under the Creative Commons Attribution 4.0 International licence, http://creativecommons.org/licenses/by/4.0/, as recorded in the arXiv licence metadata for this exact version and shown on the abstract page https://arxiv.org/abs/2512.12621v2. Full licence notice: \texttt{licenses/arxiv-2512.12621-CC-BY-4.0.txt}.\par\smallskip

\noindent\textit{Editorial scope.} Counted unit: the article's thm Thm-Fuglede (source0.tex lines 615--629). The article's complete original proof of it is delivered with it (source0.tex lines 630--707, one \texttt{proof} environment of 78 lines). No mathematics is written, repaired or re-derived: the statement and the proof are the article's own lines, in the article's own notation, and no retained line is substituted or re-flowed.  Retained uncounted as the source-local context the proof needs: lines 1621--1627; lines 1686--1691; lines 1714--1721.  Nothing was omitted from any retained span, so the package carries no omission record.  No retained line contains a citation, so the wrapper carries no bibliography and no bibliography entry.  No retained line contains a figure environment, an image-inclusion command or drawing code, so no image is delivered and the package declares no asset dependency.  Editorial formatting only, in addition: an AMS \texttt{amsart} preamble that restates the article's own declarations and macros used by the retained lines, namely the environments \texttt{thm}, \texttt{lem} on separate counters, together with the standard packages those lines need.  The retained environments print in this document as Theorem 1, Lemma 1 in order of appearance, whereas the article prints the same items with the numbering of its own full document.  The complete native source of the article as distributed in the arXiv e-print for this exact version is bundled unchanged as \texttt{sources/190959/source0.tex}.
\begin{lem}\label{Lem-appen--1}
	There exists a constant $\varepsilon_0\in(0,\frac{1}{2}]$ depending on $n$ and $R_0$, such that for $0<t\leq 2\varepsilon_0$,  we have
	\begin{equation}\label{Eq-a5}
	f_{|x-y|}\left((1+tu(x))\rho,(1+tu(x))s\right)\geq f_{|x-y|}(\rho,s)\left(1+tuG_{|x-y|}(\alpha,\rho,s)+t^2u^2D\right)
	\end{equation}
	for some function $G_{|x-y|}(\alpha,\rho,s)$ and a positive constant $D$ depending on $n,\alpha$ and $R_0$. Moreover, there exists a constant $D'$ depending on $n,\alpha$ and $R_0$, such that $\big|G_{|x-y|}(\alpha,\rho,s)\big|\leq D'$.
\end{lem}

\begin{lem}\label{Lem-appen_0}
	Under the choice of $\varepsilon_0$ in Lemma \ref{Lem-appen--1}, there exists a constant $L'$ depending on $n,\alpha$ and $R_0$, such that 
	\begin{equation}
		NL_{\alpha}(B_r)-I\leq t^2L'\sinh^{2n-\alpha} r||u||^2_{L^2(\mathbb{S}^{n-1})}.
	\end{equation}
\end{lem}

\begin{lem}\label{Lem-appen}
	%Let $f_{\theta}(\rho,s)$ be the function defined in \eqref{Eq-F5} with $\alpha\in(0,n)$, $R_0>0$ is a fixed constant, $u\in C^1(\mathbb{S}^{n-1})$ satisfiying $||u||_{C^1(\mathbb{S}^{n-1})}\leq\frac{1}{2}$. 
	Assume that $x$ and $y$ are two arbitrary points in $\mathbb{S}^{n-1}$, the values of $\rho, s$ are taken between $u(x)$ and $u(y)$ and $0<\varepsilon_0\leq\frac{1}{2}$. Then for any $\tau\in(0,2\varepsilon_0)$ and $r\in(0,R_0]$, we have 
	\begin{equation}
			\left|\frac{\partial f_{|x-y|}(r(1+\tau\rho),r(1+\tau s))}{\partial\tau}\right|\leq L''(n,\alpha,R_0)f_{|x-y|}(r,r)
	\end{equation}
for some positive constant $L''$ depending on $n,\alpha$ and $R_0$.
\end{lem}

\begin{thm}\label{Thm-Fuglede}
	For any $n\geq 2$, $0<\alpha<n$ and $R_0>0$, there exists a constant $\varepsilon_0\in(0,\frac{1}{2}]$ depending only on $n$ and $R_0$, such that if $E\subset\mathbb{H}^n$ is a nearly spherical set with $|E|_g=|B_r|_g$ for some $r\in(0,R_0]$, whose boundary is given by 
	\begin{equation}\label{Eq-Fuglede}
		\partial E=\{(x,r(1+u(x))):x\in\mathbb{S}^{n-1}\}
	\end{equation}
	under the geodesic polar coordinate, where $u:\mathbb{S}^{n-1}\to\mathbb{R}$ is a $C^1$ function, with $||u||_{C^1(\mathbb{S}^{n-1})}\leq\varepsilon_0$. Then we have the estimate 
	\begin{equation}\label{Eq-F1}
		NL_{\alpha}(B_r)-NL_{\alpha}(E)\leq L_0{\sinh}^{2n-\alpha} r\left([u]^2_{\frac{\alpha+1-n}{2}}+||u||^2_{L^2(
			\mathbb{S}^{n-1})}\right)
	\end{equation}
	for some constant $L_0>0$ depending only on $n,\alpha$ and $R_0$. Here $[u]^2_{\frac{\alpha+1-n}{2}}$ is defined by
	\begin{equation*}
		[u]^2_{\frac{\alpha+1-n}{2}}=\int_{\mathbb{S}^{n-1}}\int_{\mathbb{S}^{n-1}}\frac{|u(x)-u(y)|^2}{|x-y|^{\alpha}}\,d\sigma_x\,d\sigma_y.
	\end{equation*}
\end{thm}  

\begin{proof}
	We slightly change the notation, we assume that $|E_t|_g=|B_r|_g$, where 
	\begin{equation*}
		\partial E_t=\{(x,r(1+tu(x))):x\in\mathbb{S}^{n-1}\},
\end{equation*}
which satisfies $||u||_{C^1(\mathbb{S}^{n-1})}\leq\frac{1}{2}$ and $t\in(0,2\varepsilon_0]$. Then for any two points $\tilde{x}=(x,\rho)$ and $\tilde{y}=(y,s)$ under the geodesic polar coordinate, the distance $d_g(\tilde{x},\tilde{y})$ satisfies
\begin{equation}\label{Eq-F3}
	\cosh d_g(\tilde{x},\tilde{y})=\cosh\rho\cosh s-\sinh\rho\sinh s\cos\beta,
\end{equation}
where we have used the hyperbolic cosine law and $\beta$ is the angle between $x$ and $y$,
which satisfies
\begin{equation}\label{Eq-F4}
	\cos\beta=\frac{2-|x-y|^2}{2}.
\end{equation}
Combining \eqref{Eq-F3} with \eqref{Eq-F4}, we have
\begin{align*}
	\cosh d_g(\tilde{x},\tilde{y})&=\cosh\rho\cosh s-\sinh\rho\sinh s+\sinh\rho\sinh s(1-\cos\beta)\\
	&=\cosh(\rho-s)+\frac{\sinh\rho\sinh s}{2}|x-y|^2,
\end{align*} 
which gives
\begin{equation*}
	d_g(\tilde{x},\tilde{y})=\mathrm{arccosh}\left[\cosh(\rho-s)+\frac{\sinh\rho\sinh s}{2}|x-y|^2\right].
\end{equation*}
If we define the function $f_{\theta}(\rho,s)$ as
\begin{equation}\label{Eq-F5}
	f_{\theta}(\rho,s)=\frac{\sinh^{n-1}\rho\sinh^{n-1}s}{\mathrm{arccosh}^{\alpha}\left[\cosh(\rho-s)+\frac{\sinh\rho\sinh s}{2}{\theta}^2\right]}.
\end{equation}
Then $NL_{\alpha}(E_t)$ can be expressed by 
\begin{equation}\label{Eq-F6}
	NL_{\alpha}(E_t)=\int_{\mathbb{S}^{n-1}}\,d\sigma_x\int_{\mathbb{S}^{n-1}}\,d\sigma_y\int_0^{r(1+tu(x))}\int_0^{r(1+tu(y))}{f_{|x-y|}(\rho,s)}\,d\rho\,ds.
\end{equation}
By exploiting the identity
\begin{equation*}
	2\int_0^a\int_0^b=\int_0^a\int_0^a+\int_0^b\int_0^b-\int_a^b\int_a^b,
\end{equation*}
we deduce that 
\begin{align}
	NL_{\alpha}(E_t)=&\int_{\mathbb{S}^{n-1}}\,d\sigma_x\int_{\mathbb{S}^{n-1}}\,d\sigma_y\int_0^{r(1+tu(x))}\int_0^{r(1+tu(x))}{f_{|x-y|}(\rho,s)}\,d\rho\,ds\notag\\
	&-\frac{1}{2}\int_{\mathbb{S}^{n-1}}\,d\sigma_x\int_{\mathbb{S}^{n-1}}\,d\sigma_y\int_{r(1+tu(y))}^{r(1+tu(x))}\int_{r(1+tu(y))}^{r(1+tu(x))}{f_{|x-y|}(\rho,s)}\,d\rho\,ds\notag\\
	=:&I-II.\label{Eq-F7}
\end{align}
By taking $u=0$ in \eqref{Eq-F7}, we have
\begin{equation}\label{Eq-F8}
	NL_{\alpha}(B_r)=\int_{\mathbb{S}^{n-1}}\,d\sigma_x\int_{\mathbb{S}^{n-1}}\,d\sigma_y\int_0^{r}\int_0^{r}{f_{|x-y|}(\rho,s)}\,d\rho\,ds.
\end{equation}
Combining \eqref{Eq-F7} with \eqref{Eq-F8} gives
\begin{equation}\label{Eq-F9}
		NL_{\alpha}(B_r)-NL_{\alpha}(E_t)=NL_{\alpha}(B_r)-I+\frac{t^2 r^2}{2}h(t),
\end{equation}
where $h(t)$ is defined as 
\begin{equation}\label{Eq-F10}
	h(t)=\int_{\mathbb{S}^{n-1}}\,d\sigma_x\int_{\mathbb{S}^{n-1}}\,d\sigma_y\int_{u(y)}^{u(x)}\int_{u(y)}^{u(x)}{f_{|x-y|}(r(1+t\rho),r(1+ts))}\,d\rho\,ds.
\end{equation}
By Lemma \ref{Lem-appen_0} and \ref{Lem-appen}, there exists two constants $L'$ and $L''$ depending on $n,\alpha$ and $R_0$, such that  
\begin{align}
	&NL_{\alpha}(B_r)-I\leq t^2L' \sinh^{2n-\alpha}r||u||^2_{L^2(\mathbb{S}^{n-1})},\label{Eq-F10.2}\\
	&\left|\frac{\partial f_{|x-y|}(r(1+\tau\rho),r(1+\tau s))}{\partial\tau}\right|\leq L''f_{|x-y|}(r,r),\,\,\forall\tau\in(0,t)\label{Eq-F10.5}.
\end{align}
From \eqref{Eq-F10.5} we have
\begin{equation}\label{Eq-F11}
	|h'(\tau)|\leq L''h(0),\,\,\forall\tau\in(0,t).
\end{equation}
Next we estimate $h(0)$. Let $w=\mathrm{arccosh}\left[1+\frac{\sinh^2 r}{2}|x-y|^2\right]$, i.e. $\cosh w=1+\frac{\sinh^2 r}{2}|x-y|^2$. Since $|x-y|\leq 2$, $r\leq R_0$, there exists $L'''(R_0)$, such that $\cosh w\leq 1+\frac{L'''}{2}w^2$, which implies
\begin{equation}\label{Eq-F12}
	w\geq\frac{\sinh r}{\sqrt{L'''}}|x-y|.
\end{equation}
Combining \eqref{Eq-F10}-\eqref{Eq-F12}, we get
\begin{align}
	h(t)&\leq \left(1+tL''\right)h(0)\notag\\
	&=(1+tL'')\int_{\mathbb{S}^{n-1}}\,d\sigma_x\int_{\mathbb{S}^{n-1}}\,d\sigma_y\int_{u(y)}^{u(x)}\int_{u(y)}^{u(x)}{f_{|x-y|}(r,r)}\,d\rho\,ds\notag\\
	&\leq (1+tL'')(L''')^{\frac{\alpha}{2}}\sinh^{2n-2-\alpha} r [u]^2_{\frac{\alpha+1-n}{2}}.\label{Eq-F13}
\end{align}
Then combining \eqref{Eq-F9}, \eqref{Eq-F10.2} and \eqref{Eq-F13} gives
\begin{equation}
	NL_{\alpha}(B_r)-NL_{\alpha}(E_t)\leq t^2\sinh^{2n-\alpha} r\left[L'||u||^2_{L^2(\mathbb{S}^{n-1})}+\frac{1+tL''}{2}(L''')^{\frac{\alpha}{2}}[u]^2_{\frac{\alpha+1-n}{2}}\right],
\end{equation}
which obviously implies the announced result.
\end{proof}

\end{document}
